% 附录D 相干态路径积分
\chapter{相干态路径积分}

\section{路径积分的构造}

虚时生成泛函（见 (B.17)）为

\begin{equation*}
Z [ j, \mu ] = \operatorname{Tr} T _ {\tau} \left( \exp \left[ - \int_ {0} ^ {\beta} d \tau \, \mathcal {H} [ j ( \tau ) ] \right] \right) ,
\tag{D.1}
\end{equation*}

其中

\begin{equation*}
\mathcal {H} [ j ( \tau ) ] = \mathcal {H} _ {0} ( \mu ) - \sum_{i, \alpha} j_{i} ^ {\alpha} ( \tau ) a_{i, s} ^ {\dagger} \sigma_{ss'} ^ {\alpha} a_{i, s '} \, , \quad \tau \in [ 0, \beta ).
\tag{D.2}
\end{equation*}

源流 $j$ 与自旋算符耦合\footnote{这里我们只限于自旋源项；该形式体系可应用于任何算符的源项。}。迹公式 (C.21) 给出表示

\begin{align*}
Z & = \operatorname{Tr} T _ {\tau} \left( \exp \left[ - \int_ {0} ^ {\beta} d \tau \, \mathcal {H} [ j ( \tau ) ] \right] \right)\\
& = \int d ^ {2} \mathbf {z} _ {0} \, e ^ {- \mathbf {z} _ {0} ^ {*} \mathbf {z} _ {0}} \left\langle \mathbf {z} _ {0} \left| T \left( \exp \left[ - \int_ {0} ^ {\beta} d \tau \, H ( \tau ) \right] \right) \right| \eta \mathbf {z} _ {0} \right\rangle ,
\tag{D.3}
\end{align*}

$\mathbf{z} = (z_{1},\ldots ,z_{\mathcal{N}})$ 是在 $\mathcal{N}$ 个单粒子态基底上定义的 Bose（Fermi）相干态变量；玻色子（费米子）取 $\eta = 1$（$\eta = -1$）。

时序指数定义为

\begin{equation*}
T _ {\tau} \left( e ^ {- \int_ {0} ^ {\beta} d \tau \, \mathcal {H} ( \tau )} \right) = \lim_{\epsilon \rightarrow 0} T_{n} \prod_{n = 1} ^ {N _ {\epsilon} - 1} \left[ 1 - \epsilon \mathcal {H} ( n \epsilon ) \right] ,
\tag{D.4}
\end{equation*}

$\epsilon = \beta/N_{\epsilon}$ 是时间步长，$T_{n}$ 把因子按 $n$ 从左到右递减排序。在 (D.4) 的每两个因子之间可以插入用 Bose 或 Fermi 相干态表示的单位分解（见 (C.17)），引入 $N_{\epsilon}$ 个积分变量 $z_{\tau}$。两端等同（差一个 $\eta$ 因子）：

\begin{equation*}
\mathbf {z} _ {\beta} = \eta \mathbf {z} _ {0}.
\tag{D.5}
\end{equation*}

离散生成泛函表示为多元复（或 Grassmann）积分：

\begin{equation*}
Z ( \epsilon ) = \int \prod_{\tau = \epsilon} ^ {\beta} \left( d ^ {2} \mathbf {z} _ {\tau} \right) \exp \left[ - \sum_{\tau = \epsilon} ^ {\beta} \mathbf {z} _ {\tau} ^ {*} \left( \mathbf {z} _ {\tau} - \mathbf {z}_{\tau - \epsilon} \right) - \sum_{\tau = \epsilon} ^ {\beta} H ( \tau ) \epsilon \right] ,
\tag{D.6}
\end{equation*}

这里用 (C.19) 把 $H[\mathbf{z}^*,\mathbf{z}]$ 定义为

\begin{equation*}
\left\langle \mathbf {z} _ {\tau} \left| \mathcal {H} \right| \mathbf {z}_{\tau - \epsilon} \right\rangle = H \left[ \mathbf {z} _ {\tau} ^ {*}, \mathbf {z}_{\tau - \epsilon}, j ( \tau ) \right] e ^ {\mathbf {z} _ {\tau} ^ {*} \mathbf {z}_{\tau - \epsilon}} ,
\tag{D.7}
\end{equation*}

且 $(1 - \epsilon H) = e^{-\epsilon H} + \mathcal{O}(\epsilon)$。$\epsilon \to 0$ 的极限形式上记作路径积分

\begin{align*}
Z & = \lim_{\epsilon \to 0} Z ( \epsilon )\\
& \equiv \int \mathcal {D} ^ {2} z ( \tau ) \, \exp \left( - \int_ {0} ^ {\beta} d \tau \left\{ \mathbf {z} ^ {*} \partial_ {\tau} \mathbf {z} + H [ j ( \tau ) ] \right\} \right).
\tag{D.8}
\end{align*}

“时间导数”或“动能项”在离散表述中定义为

\begin{equation*}
\partial_ {\tau} \equiv \left( \delta_{\tau, \tau} - \delta_{\tau - \epsilon, \tau '} \right) / \epsilon .
\tag{D.9}
\end{equation*}

动能项部分地压低快变路径的贡献。

警告：由于相邻时间步的积分变量完全独立，不能假设路径 $\mathbf{z}(\tau)$ 连续或光滑！

\section{正规双线性哈密顿量}

生成泛函的一个重要子类涉及正规双线性哈密顿量：

\begin{equation*}
\int_ {0} ^ {\beta} d \tau \, H ( \tau ) = \sum_{\tau ii '} z_{i \tau} ^ {*} h_{i, i '} ( \tau ) z_{i ' \tau - \epsilon} \, \epsilon .
\tag{D.10}
\end{equation*}

正规双线性哈密顿量的生成泛函 $Z(\epsilon)$ 是多维 Gauss 积分，由 (C.12) 给出

\begin{align*}
Z ( \epsilon ) & = \int \prod_{\tau = 0} ^ {\beta - \epsilon} \left( d ^ {2} \mathbf {z} _ {\tau} \right) \exp \left[ - \epsilon \sum_{i \tau, i ' \tau '} z_{i \tau} ^ {*} G_{i \tau, i ' \tau '} ^ {- 1} z_{i ' \tau '} \right]\\
& = \left( \det G \epsilon \right) ^ {\eta}\\
& = \exp \left[ \eta \, \mathrm{Tr} \ln \left( G \epsilon \right) \right] ,
\tag{D.11}
\end{align*}

Green 函数由

\begin{equation*}
G ^ {- 1} \epsilon = \left( \hat {\partial} _ {\tau} + \hat {h} \right) \epsilon .
\tag{D.12}
\end{equation*}

给出。$G$ 可显式写成空间--时间表象中大小 $(N_{\epsilon}\mathcal{N})^2$ 的矩阵：

\begin{equation*}
G ^ {- 1} \epsilon = \left( \begin{array}{ccccc} I & 0 & 0 & \ldots & \eta \left( - I + \epsilon h ( \beta ) \right) \\ - I + \epsilon h ( 0 ) & I & 0 & \ldots & 0 \\ 0 & - I + \epsilon h ( \epsilon ) & I & 0 & \vdots \\ \vdots & \ddots & & \ddots & \ddots \\ 0 & \ldots & & - I + \epsilon h ( \beta - \epsilon ) & I \end{array} \right)
\tag{D.13}
\end{equation*}

用恒等式

\begin{equation*}
\det \left( \begin{array}{cc} A & B \\ C & D \end{array} \right) = \det D \, \det \left( A - B D ^ {- 1} C \right)
\tag{D.14}
\end{equation*}

逐次消去最后时间步的块，可把 $Z(\epsilon)$ 写成时序乘积对空间指标的行列式：

\begin{align*}
Z ( \epsilon ) & = \det \left[ I - \eta T_{n} \prod_{n = 1} ^ {N _ {\epsilon}} \left[ I - \epsilon h ( n \epsilon ) \right] \right] ^ {- \eta}\\
& \rightarrow \det \left[ I - \eta T _ {\tau} \left( \exp \left[ - \int_ {0} ^ {\beta} d \tau \, h ( \tau ) \right] \right) \right] ^ {- \eta}.
\tag{D.15}
\end{align*}

对不依赖时间的哈密顿量，可对角化 $\hat{h}_{ii'}$：

\begin{equation*}
h _ {ii '} \rightarrow \left( \epsilon_ {k} - \mu \right) \delta_{kk '} ,
\tag{D.16}
\end{equation*}

恢复无相互作用玻色子或费米子的自由能：

\begin{equation*}
F = \left\{ \begin{array}{ll} \beta^ {- 1} \sum_ {k} \ln \left( 1 - e ^ {- \beta \left( \epsilon_ {k} - \mu \right) } \right) & \text{玻色子} \\ - \beta^ {- 1} \sum_ {k} \ln \left( 1 + e ^ {- \beta \left( \epsilon_ {k} - \mu \right) } \right) & \text{费米子} \end{array} \right. ,
\tag{D.17}
\end{equation*}

此前已在 (A.24) 导出。

对有限温度的有限格子，$Z[j(\tau)]$ 可以对任意依赖时间的源项数值计算。形如 (D.15) 的公式十分有用，例如在 Hubbard 模型的量子 Monte Carlo 模拟中\footnote{见《The Hubbard Model》，第 3 章文献注释。}：$Z[Q(\tau)]$——其中 $Q(\tau)$ 是解耦四费米子相互作用的辅助场——用 (D.15) 计算，结果给出 $Q(\tau)$ 的有效 Boltzmann 权重，再用 Monte Carlo（Metropolis）算法更新。但须警惕：$Z[Q]$ 并非正定，其负号对此类模拟是严重的问题。

\section{Matsubara 表示}

为计算动力学响应函数，用 Matsubara 表示对角化 $G(\tau, \tau')$ 是有用的。以无穷时间步极限定义到 Matsubara 表示的 Fourier 变换：

\begin{align*}
X _ {\omega} & = \lim_{\epsilon \to 0} \sum_{\tau = 0} ^ {\beta - \epsilon} e ^ {- i\omega\tau} X _ {\tau} \, \epsilon\\
& = \int_ {0} ^ {\beta} d \tau \, e ^ {- i\omega\tau} X ( \tau ).
\tag{D.18}
\end{align*}

逆变换由 Matsubara 求和给出：

\begin{equation*}
X ( \tau ) = \beta^ {- 1} \sum_{n = - \infty} ^ {\infty} e ^ {i\omega_{n} \tau} X _ {\omega_ {n}} ,
\tag{D.19}
\end{equation*}

$\omega_{n}$ 是 Bose 或 Fermi 的 Matsubara 频率：

\begin{equation*}
\omega_ {n} = \left\{ \begin{array}{ll} 2 n \pi / \beta & \text{Bose} \\ ( 2n + 1 ) \pi / \beta & \text{Fermi} \end{array} \right..
\tag{D.20}
\end{equation*}

路径积分的测度可换成 Matsubara 表示：

\begin{equation*}
\mathcal {D} ^ {2} \mathbf {z} = \prod_{i, \tau} d ^ {2} z_{i, \tau} = \prod_{i, n} \frac {d ^ {2} z_{i , \omega_ {n}}}{\beta} ,
\tag{D.21}
\end{equation*}

生成泛函为

\begin{equation*}
Z = \oint \mathcal {D} ^ {2} z \, \exp \left[ - \sum_ {n} \left( - i \beta^ {- 1} \omega_ {n} \mathbf {z} _ {\omega} ^ {*} \mathbf {z} _ {\omega} \right) - \int_ {0} ^ {\beta} d \tau \, H [ j ] \right] .
\tag{D.22}
\end{equation*}

用 (D.12) 与 (D.18)，不依赖时间的双线性哈密顿量的 Green 函数的 Matsubara 表示为

\begin{equation*}
G _ {0} ( k, i\omega ) = - \frac {\delta_{nn'} \delta_{kk'}}{i\omega - h_{k}}.
\tag{D.23}
\end{equation*}

\section{Matsubara 求和}

对在无穷远处有界

\begin{equation*}
\left| g ( z ) \right| \leq \frac {c}{| z |} \, , \; | z | > R _ {0}
\tag{D.24}
\end{equation*}

的复解析函数 $g(z)$，$\tau > 0$ 的 Matsubara 求和由 Cauchy 积分公式给出：

\begin{align*}
\beta^ {- 1} \sum_ {n} e ^ {i\omega_ {n} \tau} g \left( i\omega_ {n} \right) & = - \eta \sum_{\alpha} e ^ {z_{\alpha} \tau} n \left( z_{\alpha} \right) \operatorname{Res} \left[ g \left( z_{\alpha} \right) \right]\\
& \quad - \frac {\eta}{2 \pi i} \sum_{\gamma} \int_{C _ {\gamma}} d z \, e ^ {z \tau} n ( z ) g ( z) ,
\tag{D.25}
\end{align*}

$z_{\alpha}$、$\operatorname{Res}[g(z_{\alpha})]$ 与 $C_{\gamma}$ 分别是 $g(z)$ 在复平面上的极点、留数与割线。$n(z)$ 是 Bose 或 Fermi 函数

\begin{equation*}
n ( z ) = \frac {1}{e ^ {\beta z} - \eta} ,
\tag{D.26}
\end{equation*}

它在 Matsubara 频率 $z = i\omega_{n}$ 处有单极点，留数为

\begin{equation*}
\operatorname{Res} \left[ n \left( i\omega_ {n} \right) \right] = \frac {\eta}{\beta}.
\tag{D.27}
\end{equation*}

对 $\tau > 0$，对 (D.23) 作 Matsubara 求和得虚时 Green 函数：

\begin{align*}
G _ {0} ( k, \tau ) & = \beta^ {- 1} \sum_ {n} e ^ {i\omega_ {n} \tau} G _ {0} \left( k, \omega_ {n} \right)\\
& = \eta \, e ^ {h_{k} \tau} n \left( h_{k} \right) = \eta \, e ^ {\left( \epsilon_{k} - \mu \right) \tau} n \left( \epsilon_{k} - \mu \right).
\tag{D.28}
\end{align*}

另一个有用的例子是计算自旋 $\frac{1}{2}$ 玻色子或费米子无相互作用体系的动力学自旋磁化率。本征态以格子动量 $k$ 与自旋指标 $s = \pm\frac{1}{2}$ 标记。源项写成 Matsubara 表示：

\begin{align*}
& \int_ {0} ^ {\beta} d \tau \sum_{i, \alpha} j_{i} ^ {\alpha} ( \tau ) z_{i, s} ^ {*} ( \tau ) \sigma_{ss'} ^ {\alpha} z_{i, s '} ( \tau )\\
& = - \beta^ {- 1} \sum_{\mathbf {k} n \mathbf {q} m} j_{\mathbf {q}} ^ {\alpha} \left( \nu_{m} \right) \sum_{ss'} z_{\mathbf {k} \omega_ {n}, s} ^ {*} \sigma_{ss'} ^ {\alpha} z_{\mathbf {k} + \mathbf {q} \omega_{n} + \nu_{m}, s '}\\
& \equiv \mathbf {z} ^ {\dagger} \hat {j} \mathbf {z} ,
\tag{D.29}
\end{align*}

其中

\begin{equation*}
j _ {\mathbf {q}} ^ {\alpha} ( \nu ) = \sum_ {i} \int_ {0} ^ {\beta} d \tau \, j_{i} ^ {\alpha} ( \tau ) e ^ {- i\nu\tau - i \mathbf {q} \cdot \mathbf {x} _ {i}}.
\tag{D.30}
\end{equation*}

由于选择 $j_{i}^{\alpha}(\tau)$ 在区间 $[0,\beta]$ 上周期，其 Matsubara 频率属 Bose 型：

\begin{equation*}
\nu_ {m} = 2 m \pi / \beta .
\tag{D.31}
\end{equation*}

生成泛函为

\begin{align*}
\ln Z [ j ] & = - \eta \operatorname{Tr} \ln \left[ \epsilon \left( G _ {0} ^ {- 1} - \hat {j} \right) \right]\\
& = \eta \operatorname{Tr} \ln \left( \epsilon G _ {0} \right) - \eta \sum_{n = 0} ^ {\infty} \frac {1}{n} \operatorname{Tr} \left( G _ {0} \hat {j} \right) ^ {n}.
\tag{D.32}
\end{align*}

由 (B.19) 与 (B.24)，磁化率由该展开的二阶项给出：

\begin{align*}
\chi^ {zz} ( \mathbf {q}, i\nu ) & = \frac {\partial^ {2} \ln Z}{\partial j_{\mathbf {q}} ( i\nu ) \, \partial j_{- \mathbf {q}} ( - i\nu )}\\
& = \frac {\eta}{2 \beta} \sum_{\mathbf {k}, n} G _ {0} \left( \mathbf {k}, i\omega_ {n} \right) G _ {0} \left( \mathbf {k} + \mathbf {q}, i\omega_{n} + i\nu \right)\\
& = - \frac {1}{2} \sum_{\mathbf {k}} \frac {n_{\mathbf {k}} - n_{\mathbf {k} + \mathbf {q}}}{\epsilon_{k} - \epsilon_{\mathbf {k} + \mathbf {q}} + i\nu} ,
\tag{D.33}
\end{align*}

$n_{\mathbf{k}} = n(\epsilon_{\mathbf{k}} - \mu)$。如 (B.26) 所示，解析延拓 $i\nu \to \omega + i0^{+}$ 给出实部与虚部的动力学自旋磁化率：

\begin{align*}
\operatorname{Re} \chi^ {zz} ( \mathbf {q}, \omega ) & = - \frac {1}{2} \mathcal {P} \sum_{\mathbf {k}} \frac {n_{\mathbf {k}} - n_{\mathbf {k} + \mathbf {q}}}{\epsilon_{k} - \epsilon_{\mathbf {k} + \mathbf {q}} + \omega} ,\\
\operatorname{Im} \chi^ {zz} ( \mathbf {q}, \omega ) & = \frac {\pi}{2} \sum_{\mathbf {k}} \left( n_{\mathbf {k}} - n_{\mathbf {k} + \mathbf {q}} \right) \delta \left( \epsilon_{k} - \epsilon_{\mathbf {k} + \mathbf {q}} + \omega \right).
\tag{D.34}
\end{align*}

\begin{exercises}

\item 等时 Green 函数定义为极限

\begin{equation*}
\lim_{\tau \rightarrow 0 ^ {-}} G \left( kk, \tau \right) = \left\langle a _ {k} a _ {k} ^ {\dagger} \right\rangle .
\tag{D.35}
\end{equation*}

用上述围道积分方法计算 Matsubara 求和

\begin{equation*}
G _ {0} \left( k, - \tau \right) = - \eta \beta^ {- 1} \sum_ {n} \frac {e ^ {- i\omega_ {n} \tau}}{i\omega_ {n} - \epsilon_{k}}
\tag{D.36}
\end{equation*}

并验证

\begin{equation*}
\left\langle a _ {k} a _ {k} ^ {\dagger} - \eta \, a _ {k} ^ {\dagger} a _ {k} \right\rangle = 1.
\tag{D.37}
\end{equation*}

提示：(D.36) 中的被求和项在 $i\omega_{n} \rightarrow -i\omega_{n}$ 下并非不变！

\end{exercises}

\begin{chapbib}
\item 本附录内容的推荐参考书：
  J. W. Negele and H. Orland, \textit{Quantum Many Particle Systems} (Addison-Wesley, 1988)。
\item 另见第 1、2、10 章的文献注释。
\end{chapbib}
